% !TeX root = ../paper.tex
\hypertarget{related-work}{%
\section{Related Work}\label{related-work}}

\hypertarget{generalist-vision-language-action-policies}{%
\subsection{Generalist Vision-Language-Action
Policies}\label{generalist-vision-language-action-policies}}

Generalist VLA models combine the semantic understanding capabilities of vision-language models with
robot control, improving task execution and generalization through advances in backbones, data
scale, and action modeling. The RT family demonstrates the scalability of language-conditioned
policies trained on large-scale robot data~\citep{brohan2022rt,brohan2023rt,openx2023embodiment},
while OpenVLA integrates an open VLM backbone with robot action
prediction~\citep{kim2024openvla,kim2025fine}. Building on a flow-matching action expert, the $\pi$
family combines heterogeneous robot data, high-level semantic prediction, and multimodal co-training
to improve the open-world generalization and steerability of continuous-action
policies~\citep{black2024pi,intelligence2025pi,intelligence2025pi06,intelligence2026pi07}. VLA work
in the Qwen family explores unified alignment of multi-source robot data at the representation,
action, and behavior levels, as well as the role of data scale and alignment quality in
generalization across environments, embodiments, and
instructions~\citep{wang2026qwenvla,yuan2026qwenrobotmanip}. However, task execution and
generalization performance under task-level instructions do not directly establish whether a policy
can reliably translate local commands, explicit spatial cues, and requirements updated during
execution into subsequent actions. SmartVLA complements these efforts without introducing a new
backbone architecture. Instead, it combines action-aligned data construction with staged training to
transform existing robot trajectories into semantic and spatial supervision aligned with local
control phases, learning a mapping from current instructions to local actions that supports
fine-grained instruction following, responses to spatial cues, and steering during execution.

\hypertarget{agentic-harnesses-and-callable-vla-interfaces}{%
\subsection{Agentic Harnesses and Callable VLA
Interfaces}\label{agentic-harnesses-and-callable-vla-interfaces}}

Recent agentic robotic systems organize execution as a closed loop linking high-level reasoning and
orchestration with low-level skill invocation. Works such as CaP-X~\citep{fu2026capx},
ASPIRE~\citep{lu2026aspire}, and Zetta~\citep{ding2026zetta} improve task decomposition, memory,
tool use, and failure handling through code-based policies, skill libraries and search, and runtime
critique and recovery. Stronger general-purpose multimodal models such as GPT-6
Astra~\citep{openai2026gpt6astra} further advance long-horizon reasoning, visual state
understanding, tool use, and dynamic replanning, enabling high-level harnesses to continuously
orchestrate low-level skills based on current observations and execution feedback. However, stronger
high-level reasoning and orchestration do not automatically ensure that low-level policies reliably
follow updated commands and satisfy spatial requirements. SmartVLA complements these efforts: the
high-level harness handles subgoal generation, skill selection, and plan revision, while SmartVLA
uses the current command, visual cues, and robot state to translate semantic and spatial
requirements into local actions and respond to updated instructions in subsequent calls. We compare
policies under a fixed high-level harness, replacing only the low-level VLA, to assess how low-level
instruction execution affects system performance without attributing gains from high-level planning,
prompt engineering, or tool selection to the low-level policy.

\hypertarget{steerable-and-instruction-generalizable-policies}{%
\subsection{Steerable and Instruction-Generalizable
Policies}\label{steerable-and-instruction-generalizable-policies}}

Another line of work focuses on fine-grained language, steerability, and instruction generalization.
FineVLA~\citep{hu2026finevla} uses fine-grained language supervision to provide execution-critical
details, while Steerable Policies~\citep{chen2026steerable} enables low-level policies to be invoked
through subtask commands. APT~\citep{xu2026apt} and BeTTER~\citep{xu2026better} investigate issues
such as sparse language supervision, visual shortcuts, vocabulary-motion associations, and
behavioral inertia through action-expert pretraining and spatial-layout and temporal interventions,
respectively. These works collectively emphasize that final task success rates alone cannot
establish whether the current instruction actually constrains action execution. Following this
direction, we focus on aligning semantic and spatial supervision with local actions, using data
construction and training to improve policy responsiveness to fine-grained instructions, explicit
spatial cues, and instruction updates during execution.

